%%%%%%%%%%%%%%%%%%%%%%%%%%%%%%%%%%%%%%%%%%%%%%%%%%%%%%%%%%%%%%%%%%%%%%%
\begin{comment}
\textblockcolour{pink}
\begin{textblock}{4}(16, 1)
    \noindent
    【2】2枚目は「内表紙」
\end{textblock}

\textblockcolour{pink}
\begin{textblock}{8}(0.5, 6.5)
    【6,10】タイトルが表紙・概要と同じか
\end{textblock}

\textblockcolour{pink}
\begin{textblock}{4}(0.5, 12.5)
    【6,10】指導教員名
\end{textblock}

\textblockcolour{lime}
\begin{textblock}{4}(13, 11.2)
    ↓職位の有無は任意
\end{textblock}

\textblockcolour{pink}
\begin{textblock}{5}(15.3, 13.5)
    【9】↑紙印刷の頃にあった

    押印用のセルを作らない
\end{textblock}

\textblockcolour{pink}
\begin{textblock}{8}(11, 15.8)
    \noindent
    【6,13】日付が提出期間内かを確認する 

    \noindent
    月日に1桁の数値を書くときは0埋めをしない
\end{textblock}


\begin{comment19}
    \textblockcolour{pink}
    \begin{textblock}{5}(0.5, 21.3)
        \noindent
        【6,10,11】
        
        \noindent
        学部・学籍番号・氏名
    \end{textblock}
\end{comment19}

\begin{comment20}
    \textblockcolour{pink}
    \begin{textblock}{5}(0.5, 21.3)
        \noindent
        【6,10,11】
        
        \noindent
        学部・専攻・学籍番号・氏名
    \end{textblock}
\end{comment20}

\textblockcolour{pink}
\begin{textblock}{6}(14.5, 23)
    【12】学籍番号のA,B,Cは大文字
\end{textblock}

\begin{textblock}{7}(11, 28)
    ←内表紙にページ番号をつけない
\end{textblock}
\end{comment}
%%%%%%%%%%%%%%%%%%%%%%%%%%%%%%%%%%%%%%%%%%%%%%%%%%%%%%%%%%%%%%%%%%%%%%%